% !TeX root = ../main.tex
\courseunit{3}

\begin{problem}{1}
Consider the model of Section 3.2 with \(\theta<1\).
\begin{enumerate}[label=(\alph*)]
    \item On the balanced growth path, \(\dot{A}=g_A^*A\left(t\right)\), where \(g_A^*\) is the balanced-growth-path value of \(g_A\). Use this fact and equation (3.6) to derive an expression for \(A\left(t\right)\) on the balanced growth path in terms of \(B\), \(a_L\), \(\gamma\), \(\theta\), and \(L\left(t\right)\).
    \item Use your answer to part (a) and the production function, (3.5), to obtain an expression for \(Y\left(t\right)\) on the balanced growth path. Find the value of \(a_L\) that maximizes output on the balanced growth path.
\end{enumerate}
\end{problem}

\begin{problem}{2}
Consider two economies (indexed by \(i=1,2\)) described by \(Y_i\left(t\right)=K_i\left(t\right)^\theta\) and \(\dot{K}_i\left(t\right)=s_iY_i\left(t\right)\), where \(\theta>1\). Suppose that the two economies have the same initial value of \(K\), but that \(s_1>s_2\). Show that \(Y_1/Y_2\) is continually rising.
\end{problem}

\begin{problem}{3}
Consider the economy analyzed in Section 3.3. Assume that \(\theta+\beta<1\) and \(n>0\), and that the economy is on its balanced growth path. Describe how each of the following changes affects the \(\dot{g}_A=0\) and \(\dot{g}_K=0\) lines and the position of the economy in \(\left(g_A,g_K\right)\) space at the moment of the change:
\begin{enumerate}[label=(\alph*)]
    \item An increase in \(n\).
    \item An increase in \(a_K\).
    \item An increase in \(\theta\).
\end{enumerate}
\end{problem}

\begin{problem}{4}
Consider the economy described in Section 3.3, and assume \(\beta+\theta<1\) and \(n>0\). Suppose the economy is initially on its balanced growth path, and that there is a permanent increase in \(s\).
\begin{enumerate}[label=(\alph*)]
    \item How, if at all, does the change affect the \(\dot{g}_A=0\) and \(\dot{g}_K=0\) lines? How, if at all, does it affect the location of the economy in \(\left(g_A,g_K\right)\) space at the time of the change?
    \item What are the dynamics of \(g_A\) and \(g_K\) after the increase in \(s\)? Sketch the path of log output per worker.
    \item Intuitively, how does the effect of the increase in \(s\) compare with its effect in the Solow model?
\end{enumerate}
\end{problem}

\begin{problem}{5}
Consider the model of Section 3.3 with \(\beta+\theta=1\) and \(n=0\).
\begin{enumerate}[label=(\alph*)]
    \item Using (3.14) and (3.16), find the value that \(A/K\) must have for \(g_K\) and \(g_A\) to be equal.
    \item Using your result in part (a), find the growth rate of \(A\) and \(K\) when \(g_K=g_A\).
    \item How does an increase in \(s\) affect the long-run growth rate of the economy?
    \item What value of \(a_K\) maximizes the long-run growth rate of the economy? Intuitively, why is this value not increasing in \(\beta\), the importance of capital in the R\&D sector?
\end{enumerate}
\end{problem}

\begin{problem}{6}
Consider the model of Section 3.3 with \(\beta+\theta>1\) and \(n>0\).
\begin{enumerate}[label=(\alph*)]
    \item Draw the phase diagram for this case.
    \item Show that regardless of the economy's initial conditions, eventually the growth rates of \(A\) and \(K\) (and hence the growth rate of \(Y\)) are increasing continually.
    \item Repeat parts (a) and (b) for the case of \(\beta+\theta=1\), \(n>0\).
\end{enumerate}
\end{problem}

\begin{problem}{7}[\textbf{Learning-by-doing.}]
Suppose that output is given by equation (3.22), \(Y\left(t\right)=K\left(t\right)^\alpha\left[A\left(t\right)L\left(t\right)\right]^{1-\alpha}\); that \(L\) is constant and equal to 1; that \(\dot{K}\left(t\right)=sY\left(t\right)\); and that knowledge accumulation occurs as a side effect of goods production: \(\dot{A}\left(t\right)=BY\left(t\right)\).
\begin{enumerate}[label=(\alph*)]
    \item Find expressions for \(g_A\left(t\right)\) and \(g_K\left(t\right)\) in terms of \(A\left(t\right)\), \(K\left(t\right)\), and the parameters.
    \item Sketch the \(\dot{g}_A=0\) and \(\dot{g}_K=0\) lines in \(\left(g_A,g_K\right)\) space.
    \item Does the economy converge to a balanced growth path? If so, what are the growth rates of \(K\), \(A\), and \(Y\) on the balanced growth path?
    \item How does an increase in \(s\) affect long-run growth?
\end{enumerate}
\end{problem}

\begin{problem}{8}
Consider the model of Section 3.5. Suppose, however, that households have constant-relative-risk-aversion utility with a coefficient of relative risk aversion of \(\theta\). Find the equilibrium level of labor in the R\&D sector, \(L_A\).
\end{problem}

\begin{problem}{9}
Suppose that policymakers, realizing that monopoly power creates distortions, put controls on the prices that patent-holders in the Romer model can charge for the inputs embodying their ideas. Specifically, suppose they require patent-holders to charge \(\delta w\left(t\right)/\phi\), where \(\delta\) satisfies \(\phi\leq\delta\leq1\).
\begin{enumerate}[label=(\alph*)]
    \item What is the equilibrium growth rate of the economy as a function of \(\delta\) and the other parameters of the model? Does a reduction in \(\delta\) increase, decrease, or have no effect on the equilibrium growth rate, or is it not possible to tell?
    \item Explain intuitively why setting \(\delta=\phi\), thereby requiring patent-holders to charge marginal cost and so eliminating the monopoly distortion, does not maximize social welfare.
\end{enumerate}
\end{problem}

\begin{problem}{10}
\begin{enumerate}[label=(\alph*)]
    \item Show that (3.48) follows from (3.47).
    \item Derive (3.49).
\end{enumerate}
\end{problem}

\begin{problem}{11}[\textbf{The balanced growth path of a semi-endogenous version of the Romer model.}]
\problemsource{\parencite{jones1995rd}}
Consider the model of Section 3.5 with two changes. First, existing knowledge contributes less than proportionally to the production of new knowledge, as in Case 1 of the model of Section 3.2: \(\dot{A}\left(t\right)=BL_A\left(t\right)A\left(t\right)^\theta\), \(\theta<1\). Second, population is growing at rate \(n\) rather than constant: \(L\left(t\right)=L\left(0\right)e^{nt}\), \(n>0\). (Consistent with this, assume that utility is given by equation [2.2] with \(u\left(\,\cdot\,\right)\) logarithmic.) The analysis of Section 3.2 implies that such a model will exhibit transition dynamics rather than always immediately being on a balanced growth path. This problem therefore asks you not to analyze the full dynamics of the model, but to focus on its properties when it is on a balanced growth path. Specifically, it looks at situations where the fraction of the labor force engaged in R\&D, \(a_L\), is constant, and all variables of the model are growing at constant rates.
\par \remarkaccent{Hint:} You are welcome to assume rather than derive that on a balanced growth path, the wage and consumption per person grow at the same rate as \(Y/L\).
\begin{enumerate}[label=(\alph*)]
    \item What are balanced-growth-path values of the growth rates of \(Y\) and \(Y/L\) and of \(r\) as functions of the balanced-growth-path value of the growth rate of \(A\) and parameters of the model?
    \item On the balanced growth path, what is the present value of the profits from the discovery of a new idea at time \(t\) as a function of \(L\left(t\right)\), \(A\left(t\right)\), \(w\left(t\right)\), and exogenous parameters?
    \item What is \(\dot{A}\left(t\right)/A\left(t\right)\) on the balanced growth path as a function of \(a_L\) and exogenous parameters? What is \(L\left(t\right)A\left(t\right)^{\theta-1}\) on the balanced growth path?
    \par \remarkaccent{Hint:} Consider equations [3.6] and [3.7] in the case of \(\gamma=1\).
    \item What is \(a_L\) on the balanced growth path?
    \item Discuss how changes in each of \(\rho\), \(B\), \(\phi\), \(n\), and \(\theta\) affect the balanced-growth-path value of \(a_L\). In the cases of \(\rho\), \(B\), and \(\phi\), are the effects in the same direction as the effects on \(L_A\) in the model of Section 3.5?
\end{enumerate}
\end{problem}

\begin{problem}{12}[\textbf{Learning-by-doing with microeconomic foundations.}]
Consider a variant of the model in equations (3.22)--(3.25). Suppose firm \(i\)'s output is \(Y_i\left(t\right)=K_i\left(t\right)^\alpha\left[A\left(t\right)L_i\left(t\right)\right]^{1-\alpha}\), and that \(A\left(t\right)=BK\left(t\right)\). Here \(K_i\) and \(L_i\) are the amounts of capital and labor used by firm \(i\), and \(K\) is the aggregate capital stock. Capital and labor earn their \emph{private} marginal products. As in the model of Section 3.5, the economy is populated by infinitely lived households that own the economy's initial capital stock. The utility of the representative household takes the constant-relative-risk-aversion form in equations (2.2)--(2.3). Population growth is zero.
\begin{enumerate}[label=(\alph*)]
    \item
    \begin{enumerate}[label=(\roman*)]
        \item What are the private marginal products of capital and labor at firm \(i\) as functions of \(K_i\left(t\right)\), \(L_i\left(t\right)\), \(K\left(t\right)\), and the parameters of the model?
        \item Explain why the capital-labor ratio must be the same at all firms, so
        \begin{align*}
            \dfrac{K_i\left(t\right)}{L_i\left(t\right)}=\dfrac{K\left(t\right)}{L\left(t\right)}
        \end{align*}
        \item What are \(w\left(t\right)\) and \(r\left(t\right)\) as functions of \(K\left(t\right)\), \(L\), and the parameters of the model?
    \end{enumerate}
    \item What must the growth rate of consumption be in equilibrium? Assume for simplicity that the parameter values are such that the growth rate is strictly positive and less than the interest rate. Sketch an explanation of why the equilibrium growth rate of output equals the equilibrium growth rate of consumption.
    \par \remarkaccent{Hint:} Consider equation [2.22].
    \item Describe how long-run growth is affected by:
    \begin{enumerate}[label=(\roman*)]
        \item A rise in \(B\).
        \item A rise in \(\rho\).
        \item A rise in \(L\).
    \end{enumerate}
    \item Is the equilibrium growth rate more than, less than, or equal to the socially optimal rate, or is it not possible to tell?
\end{enumerate}
\end{problem}

\begin{problem}{13}
\problemsource{\parencite{rebelo1991long-run}}
Assume that there are two sectors, one producing consumption goods and one producing capital goods, and two factors of production: capital and land. Capital is used in both sectors, but land is used only in producing consumption goods. Specifically, the production functions are \(C\left(t\right)=K_C\left(t\right)^\alpha T^{1-\alpha}\) and \(\dot{K}\left(t\right)=BK_K\left(t\right)\), where \(K_C\) and \(K_K\) are the amounts of capital used in the two sectors (so \(K_C\left(t\right)+K_K\left(t\right)=K\left(t\right)\)) and \(T\) is the amount of land, and \(0<\alpha<1\) and \(B>0\). Factors are paid their marginal products, and capital can move freely between the two sectors. \(T\) is normalized to 1 for simplicity.
\begin{enumerate}[label=(\alph*)]
    \item Let \(P_K\left(t\right)\) denote the price of capital goods relative to consumption goods at time \(t\). Use the fact that the earnings of capital in units of consumption goods in the two sectors must be equal to derive a condition relating \(P_K\left(t\right)\), \(K_C\left(t\right)\), and the parameters \(\alpha\) and \(B\). If \(K_C\) is growing at rate \(g_K\left(t\right)\), at what rate must \(P_K\) be growing (or falling)? Let \(g_P\left(t\right)\) denote this growth rate.
    \item The real interest rate in terms of consumption is \(B+g_P\left(t\right)\).\footnote{To see this, note that capital in the investment sector produces new capital at rate \(B\) and changes in value relative to the consumption good at rate \(g_P\). (Because the return to capital is the same in the two sectors, the same must be true of capital in the consumption sector.)} Thus, assuming that households have our standard utility function, (2.22)--(2.23), the growth rate of consumption must be \(\left(B+g_P-\rho\right)/\theta\equiv g_C\). Assume \(\rho<B\).
    \begin{enumerate}[label=(\roman*)]
        \item Use your results in part (a) to express \(g_C\left(t\right)\) in terms of \(g_K\left(t\right)\) rather than \(g_P\left(t\right)\).
        \item Given the production function for consumption goods, at what rate must \(K_C\) be growing for \(C\) to be growing at rate \(g_C\left(t\right)\)?
        \item Combine your answers to (i) and (ii) to solve for \(g_K\left(t\right)\) and \(g_C\left(t\right)\) in terms of the underlying parameters.
    \end{enumerate}
    \item Suppose that investment income is taxed at rate \(\tau\), so that the real interest rate households face is \(\left(1-\tau\right)\left(B+g_P\right)\). How, if at all, does \(\tau\) affect the equilibrium growth rate of consumption?
\end{enumerate}
\end{problem}

\begin{problem}{14}[\textbf{Delays in the transmission of knowledge to poor countries.}]
\begin{enumerate}[label=(\alph*)]
    \item Assume that the world consists of two regions, the North and the South. The North is described by \(Y_N\left(t\right)=A_N\left(t\right)\left(1-a_L\right)L_N\) and \(\dot{A}_N\left(t\right)=a_LL_NA_N\left(t\right)\). The South does not do R\&D but simply uses the technology developed in the North; however, the technology used in the South lags the North's by \(\tau\) years. Thus \(Y_S\left(t\right)=A_S\left(t\right)L_S\) and \(A_S\left(t\right)=A_N\left(t-\tau\right)\). If the growth rate of output per worker in the North is 3 percent per year, and if \(a_L\) is close to 0, what must \(\tau\) be for output per worker in the North to exceed that in the South by a factor of 10?
    \item Suppose instead that both the North and the South are described by the Solow model: \(y_i\left(t\right)=f\left(k_i\left(t\right)\right)\), where 
    \begin{align*}
        y_i\left(t\right) \equiv \dfrac{Y_i\left(t\right)}{L_i\left(t\right)} \quad \text{and} \quad k_i\left(t\right)\equiv \dfrac{K_i\left(t\right)}{L_i\left(t\right)}\,,\quad i=N,S
    \end{align*}
    As in the Solow model, assume \(\dot{K}_i\left(t\right)=sY_i\left(t\right)-\delta K_i\left(t\right)\) and \(\dot{L}_i\left(t\right)=nL_i\left(t\right)\); the two countries are assumed to have the same saving rates and rates of population growth. Finally, \(\dot{A}_N\left(t\right)=gA_N\left(t\right)\) and \(A_S\left(t\right)=A_N\left(t-\tau\right)\).
    \begin{enumerate}[label=(\roman*)]
        \item Show that the value of \(k\) on the balanced growth path, \(k^*\), is the same for the two countries.
        \item Does introducing capital change the answer to part (a)? Explain. (Continue to assume \(g=3\%\).)
    \end{enumerate}
\end{enumerate}
\end{problem}

\begin{problem}{15}
Which of the following possible regression results concerning the elasticity of long-run output with respect to the saving rate would provide the best evidence that differences in saving rates are not important to cross-country income differences? (1) A point estimate of 5 with a standard error of 2; (2) a point estimate of 0.1 with a standard error of 0.01; (3) a point estimate of 0.001 with a standard error of 5; (4) a point estimate of \(-2\) with a standard error of 5. Explain your answer.
\par \remarkaccent{Hint:} See the discussion of confidence intervals versus \(t\)-statistics in Section 3.6.
\end{problem}
